\daytitle{21}{磁场与安培力}

\begin{learningbox}{本节导读}
前面几节学习了电场和电路：电荷受电场力，电流表示电荷定向移动。本节进入磁场，先认识磁感应强度 $B$、磁场方向和磁感线，再学习通电导线在磁场中受到的安培力。磁场部分最重要的不是记公式，而是把“方向判断”和“受力分析”合在一起。
\end{learningbox}

\textbf{本节读题路线：} 先分清题目问的是“电流产生磁场”还是“通电导线受力”：前者用右手螺旋定则，后者先判断外磁场方向和电流方向，再用左手定则确定安培力方向。计算大小前确认有效导线长度是否在磁场中、是否与磁场垂直。若导体棒处于平衡或运动状态，把安培力和重力、弹力、摩擦力一起画进受力图。

\section*{一、学习目标}
\begin{enumerate}
  \item 能举例说明磁场能对磁体、电流和运动电荷产生作用。
  \item 能用磁感应强度 $B$ 描述磁场强弱和方向，并写出单位特斯拉（$\mathrm{T}$）。
  \item 会判断直线电流、环形电流和通电螺线管周围磁场方向。
  \item 能用磁畴和电流磁场图像说明磁化、铁磁性材料和电磁铁的基本物理图像。
  \item 能根据导线长度、电流方向和磁场夹角使用安培力公式 $F = BIL\sin\theta$，明确公式的适用条件。
  \item 会用左手定则判断安培力方向，并能把安培力放入受力图中分析平衡与运动。
\end{enumerate}
\section*{二、初高中衔接}

初中已经学过磁体有南北极、同名磁极相斥、异名磁极相吸，也知道通电导线周围存在磁场。高中要把这些现象进一步量化：
\begin{itemize}
  \item 用磁感应强度 $B$ 描述磁场强弱和方向；
  \item 用磁感线描述磁场分布，但磁感线不是实际存在的线；
  \item 用安培力描述磁场对电流的作用；
  \item 把安培力纳入力学受力分析，和重力、弹力、摩擦力一起处理。
\end{itemize}

\textbf{科学发展视角：} 磁场和安培力是电动机、扬声器、磁悬浮和电动汽车驱动系统的基础。现代设备并不是只“有磁场”就能工作，而是要控制电流方向、磁场方向和导线位置，让安培力按需要的方向输出。学本节时要把应用都翻译回三个量：$I$、$B$ 和受力方向。

\section*{三、核心概念}

\begin{conceptbox}{磁场与磁感线}

\textbf{磁场：} 磁体、电流和运动电荷周围存在磁场。磁场的基本性质是对放入其中的磁体、电流或运动电荷产生力的作用。

\textbf{磁场方向：} 规定小磁针静止时北极（$N$ 极）所指方向为该点磁场方向。

\textbf{磁感线：} 用来形象表示磁场方向和强弱的曲线。
\begin{itemize}
  \item 磁感线上某点的切线方向表示该点磁场方向。
  \item 磁感线疏密表示磁场强弱：越密，磁场越强。
  \item 磁感线是闭合曲线，在磁体外部从 $N$ 极指向 $S$ 极，在磁体内部从 $S$ 极回到 $N$ 极。
\end{itemize}

\begin{center}
% TikZ 图：条形磁体磁感线
% Day 21：磁场与安培力
\begin{tikzpicture}[scale=0.95,>=Latex]
  \draw[thick,fill=red!9] (-2.25,-0.38) rectangle (0,0.38);
  \draw[thick,fill=blue!9] (0,-0.38) rectangle (2.25,0.38);
  \draw[thick] (0,-0.38) -- (0,0.38);
  \node[font=\large] at (-1.13,0) {$N$};
  \node[font=\large] at (1.13,0) {$S$};

  \draw[blue,thick,->] (-2.22,0.30) .. controls (-2.60,1.35) and (2.60,1.35) .. (2.22,0.30);
  \draw[blue,thick,->] (-2.22,0.15) .. controls (-2.32,0.78) and (2.32,0.78) .. (2.22,0.15);
  \draw[blue,thick,->] (-2.22,-0.15) .. controls (-2.32,-0.78) and (2.32,-0.78) .. (2.22,-0.15);
  \draw[blue,thick,->] (-2.22,-0.30) .. controls (-2.60,-1.35) and (2.60,-1.35) .. (2.22,-0.30);

  \draw[blue,thick,->] (-2.05,0.38) .. controls (-1.35,0.60) and (1.35,0.60) .. (2.05,0.38);
  \draw[blue,thick,->] (-2.05,-0.38) .. controls (-1.35,-0.60) and (1.35,-0.60) .. (2.05,-0.38);
  \draw[blue!55,dashed,thick,->] (1.65,0) -- (-1.65,0);

  \node[blue] at (0,1.62) {外部磁感线：从 $N$ 极出发，进入 $S$ 极};
\end{tikzpicture}


\figcaption{条形磁体外部磁感线：外部从 $N$ 极指向 $S$ 极}\label{fig:magnetic-field-lines}
\end{center}

\end{conceptbox}

\begin{conceptbox}{地磁场与指南针}

地球本身也有磁场，指南针能指示方向，正是因为小磁针会沿当地磁场方向转动。\textbf{方向检查：} 指南针的 $N$ 极指向地理北方，说明地理北极附近在磁性上相当于\textbf{地磁 $S$ 极}，否则不会吸引指南针的 $N$ 极。

\begin{center}
\begin{tikzpicture}[scale=0.92,>=Stealth]
  \tikzset{
    field/.style={blue,thick,->},
    axis/.style={gray,dashed},
    note/.style={font=\small,align=center}
  }

  \shade[ball color=blue!12] (0,0) circle (1.25);
  \draw[thick] (0,0) circle (1.25);
  \draw[axis] (0,-1.62) -- (0,1.62);
  \node[font=\small] at (0,1.88) {地理北极附近};
  \node[font=\small] at (0,-1.88) {地理南极附近};

  \node[font=\small,red] at (1.15,1.12) {磁 $S$};
  \node[font=\small,red] at (1.15,-1.12) {磁 $N$};

  \draw[field] (0,-1.30) .. controls (-2.05,-0.70) and (-2.05,0.70) .. (0,1.30);
  \draw[field] (0,-1.30) .. controls (2.05,-0.70) and (2.05,0.70) .. (0,1.30);
  \draw[field] (0,-1.18) .. controls (-1.30,-0.52) and (-1.30,0.52) .. (0,1.18);
  \draw[field] (0,-1.18) .. controls (1.30,-0.52) and (1.30,0.52) .. (0,1.18);

  \draw[thick,orange!90!black,->] (-2.45,0.10) -- (-1.72,0.55);
  \draw[thick,orange!90!black] (-2.45,0.10) -- (-2.86,-0.16);
  \node[font=\scriptsize,orange!90!black] at (-3.08,-0.32) {指南针 $N$ 极};
  \node[font=\scriptsize,blue] at (-2.04,0.90) {沿当地磁场方向};

  \node[note] at (0,-2.42)
    {外部磁感线从磁 $N$ 极出发，进入磁 $S$ 极};
\end{tikzpicture}


\figcaption{地磁场模型：指南针 $N$ 极沿当地磁场方向，指向地理北方附近}\label{fig:earth-magnetic-field}
\end{center}

\textbf{认知提醒：} 地磁场并不完全等同于一个理想条形磁体，磁轴和地理自转轴也不完全重合。高中阶段先用“近似条形磁体”模型理解方向，再在地理和现代导航中认识磁偏角等实际修正。

\end{conceptbox}

\begin{conceptbox}{磁感应强度 $B$}

磁感应强度 $B$ 是描述磁场强弱和方向的物理量。

在匀强磁场中，当一段长度为 $L$、电流为 $I$ 的直导线垂直于磁场放置时，受到的安培力为 $F$，定义
\[
\boxed{B = \frac{F}{IL}}
\]

\begin{itemize}
  \item 单位：特斯拉（$\mathrm{T}$），$1\ \mathrm{T} = 1\ \mathrm{N/(A\cdot m)}$。
  \item $B$ 是矢量，方向就是该点磁场方向。
  \item 定义式 $B=\dfrac{F}{IL}$ 要求导线与磁场垂直；一般情形要用安培力公式。
\end{itemize}

\begin{center}
% TikZ 图：垂直纸面的磁场符号
% Day 21：磁场与安培力
\begin{tikzpicture}[scale=0.95,>=Latex]
  \node at (0,1.6) {垂直纸面向里};
  \foreach \x in {-1.2,-0.4,0.4,1.2} {
    \foreach \y in {0.2,0.9} {
      \draw[blue,thick] (\x-0.11,\y-0.11) -- (\x+0.11,\y+0.11);
      \draw[blue,thick] (\x-0.11,\y+0.11) -- (\x+0.11,\y-0.11);
    }
  }
  \node at (4.4,1.6) {垂直纸面向外};
  \foreach \x in {3.2,4.0,4.8,5.6} {
    \foreach \y in {0.2,0.9} {
      \draw[blue,thick] (\x,\y) circle (0.13);
      \fill[blue] (\x,\y) circle (0.035);
    }
  }
\end{tikzpicture}


\figcaption{垂直纸面的磁场符号：叉号表示垂直纸面向里，点号表示垂直纸面向外}\label{fig:magnetic-field-symbols}
\end{center}

\end{conceptbox}

\begin{conceptbox}{电流周围的磁场}

电流会产生磁场。高中阶段常见三类图像：
\begin{itemize}
  \item \textbf{直线电流：} 磁感线是以导线为中心的同心圆。
  \item \textbf{环形电流：} 一侧相当于 $N$ 极，另一侧相当于 $S$ 极，中心附近磁场近似沿轴线方向。
  \item \textbf{通电螺线管：} 外部磁场类似条形磁体。
\end{itemize}

\begin{center}
% TikZ 图：直线电流周围磁场
% Day 21：磁场与安培力
\begin{tikzpicture}[scale=0.92,>=Latex]
  \begin{scope}[shift={(-2.2,0)}]
    \draw[thick] (0,0) circle (0.22);
    \fill[red] (0,0) circle (0.055);
    \node[red] at (0,0.48) {$I$ 向外};
    \draw[blue,thick,->] (0.88,0) arc (0:270:0.88);
    \draw[blue,thick,->] (1.28,0) arc (0:270:1.28);
    \node[blue] at (-1.23,1.12) {$B$};
    \node[font=\small] at (0,-1.68) {拇指指出纸面};
  \end{scope}

  \begin{scope}[shift={(2.2,0)}]
    \draw[thick] (0,0) circle (0.22);
    \draw[red,thick] (-0.12,-0.12) -- (0.12,0.12);
    \draw[red,thick] (-0.12,0.12) -- (0.12,-0.12);
    \node[red] at (0,0.48) {$I$ 向里};
    \draw[blue,thick,->] (0.88,0) arc (0:-270:0.88);
    \draw[blue,thick,->] (1.28,0) arc (0:-270:1.28);
    \node[blue] at (1.23,1.12) {$B$};
    \node[font=\small] at (0,-1.68) {拇指指入纸面};
  \end{scope}

  \node[font=\small,align=center] at (0,-2.18) {俯视导线截面：右手四指弯曲方向就是磁场方向};
\end{tikzpicture}


\figcaption{直线电流的右手螺旋定则：拇指指电流方向，四指环绕方向为磁场方向}\label{fig:current-magnetic-field}
\end{center}

\begin{center}
% TikZ 图：环形电流与通电螺线管的磁场方向
% Day 21：磁场与安培力
\begin{tikzpicture}[scale=0.88,>=Latex,line cap=round,line join=round]
  \path[use as bounding box] (-5.10,-2.05) rectangle (5.10,2.05);

  % 左图：正视环形电流，逆时针电流在中心产生向外的磁场。
  \begin{scope}[shift={(-2.65,0)}]
    \node[font=\small] at (0,1.70) {环形电流：此面为 $N$ 极};
    \draw[red,very thick] (0,0) circle (1.12);
    \draw[red,very thick,->] (1.12,0) arc (0:72:1.12);
    \draw[red,very thick,->] (-1.12,0) arc (180:252:1.12);
    \node[red,font=\small] at (0,-1.45) {$I$ 沿逆时针方向};

    \draw[blue,thick] (0,0) circle (0.18);
    \fill[blue] (0,0) circle (0.052);
    \node[blue,font=\small] at (0,0.42) {$B$ 向外};
  \end{scope}

  % 右图：从左端看电流逆时针，因此左端为 N 极，管内磁场向左。
  \begin{scope}[shift={(2.55,0)}]
    \node[font=\small] at (0,1.70) {通电螺线管};
    \draw[thick] (-1.65,-0.52) rectangle (1.65,0.52);
    \foreach \x in {-1.45,-1.10,-0.75,-0.40,-0.05,0.30,0.65,1.00,1.35} {
      \draw[red,thick] (\x,-0.68)
        .. controls (\x+0.20,-0.24) and (\x+0.20,0.24) .. (\x,0.68);
    }

    \node[font=\small] at (-1.93,0) {$N$};
    \node[font=\small] at (1.93,0) {$S$};
    \draw[blue,very thick,->] (1.35,0) -- (-1.35,0);
    \node[blue,font=\small] at (0,0.30) {管内 $B$};

    \draw[blue,thick,->] (-1.55,0.82)
      .. controls (-0.45,1.35) and (0.45,1.35) .. (1.55,0.82);
    \draw[blue,thick,->] (-1.55,-0.82)
      .. controls (-0.45,-1.35) and (0.45,-1.35) .. (1.55,-0.82);

    \node[red,font=\scriptsize,align=center] at (0,-1.58)
      {从左端看电流逆时针\\四指沿电流，拇指指向左端};
  \end{scope}
\end{tikzpicture}


\figcaption{环形电流与通电螺线管的右手螺旋定则：四指沿电流方向，拇指指向轴线上磁场方向}\label{fig:loop-solenoid-field}
\end{center}

\textbf{方向规则要分清：} 对直线电流，拇指沿电流方向，四指环绕方向就是磁场方向；对环形电流和通电螺线管，四指沿电流环绕方向，拇指所指方向是中心轴线附近的磁场方向，也对应螺线管的 $N$ 极方向。反过来，从某一端看，若线圈电流为逆时针，该端为 $N$ 极；若为顺时针，该端为 $S$ 极。

\textbf{认知提醒：} 判断电流产生的磁场用右手；判断通电导线受安培力用左手。两种定则解决的问题不同，不要混用。

\end{conceptbox}

\begin{conceptbox}{磁畴、磁化与电磁铁}

铁、钴、镍等材料容易被磁化，称为铁磁性材料。铁磁性材料内部可看成含有许多微小的磁性区域，这些区域称为\textbf{磁畴}。未磁化时，各磁畴的磁化方向较为杂乱，宏观磁性很弱；在外磁场作用下，较多磁畴趋向同一方向，材料便表现出较强磁性，这一过程称为\textbf{磁化}。

把软铁芯插入通电螺线管后，软铁芯被磁化，使螺线管的磁场明显增强，这就是电磁铁的基本原理。

\begin{center}
\begin{minipage}{0.98\linewidth}
\centering
% TikZ 图：磁化与电磁铁
% Day 21：磁场与安培力
\begin{tikzpicture}[scale=0.88,>=Latex,line cap=round,line join=round]
  \path[use as bounding box] (-5.15,-2.10) rectangle (5.15,2.15);

  \begin{scope}[shift={(-2.55,0)}]
    \node[font=\small] at (0,1.78) {空心螺线管};
    \draw[thick] (-1.55,-0.42) rectangle (1.55,0.42);
    \foreach \x in {-1.30,-0.95,-0.60,-0.25,0.10,0.45,0.80,1.15} {
      \draw[red,thick] (\x,-0.58) .. controls (\x+0.18,-0.18) and (\x+0.18,0.18) .. (\x,0.58);
    }
    \draw[red,thick,->] (-1.75,-0.58) -- (-1.40,-0.58);
    \node[red,font=\small] at (-1.95,-0.78) {$I$};

    \foreach \y in {-0.22,0.22} {
      \draw[blue,thick,->] (-1.85,\y) -- (1.85,\y);
    }
    \draw[blue,thick,->] (-1.80,0.86) .. controls (-0.65,1.18) and (0.65,1.18) .. (1.80,0.86);
    \draw[blue,thick,->] (1.80,-0.86) .. controls (0.65,-1.18) and (-0.65,-1.18) .. (-1.80,-0.86);
    \node[blue,font=\small] at (1.95,0.36) {$B$};
    \node[font=\small] at (0,-1.58) {电流产生磁场};
  \end{scope}

  \begin{scope}[shift={(2.55,0)}]
    \node[font=\small] at (0,1.78) {插入软铁芯：电磁铁};
    \draw[thick,fill=gray!25] (-1.62,-0.35) rectangle (1.62,0.35);
    \foreach \x in {-1.38,-1.03,-0.68,-0.33,0.02,0.37,0.72,1.07} {
      \draw[red,thick] (\x,-0.58) .. controls (\x+0.18,-0.18) and (\x+0.18,0.18) .. (\x,0.58);
    }
    \draw[red,thick,->] (-1.82,-0.58) -- (-1.47,-0.58);
    \node[red,font=\small] at (-2.03,-0.78) {$I$};

    \foreach \y in {-0.35,0,0.35} {
      \draw[blue,very thick,->] (-1.98,\y) -- (1.98,\y);
    }
    \draw[blue,very thick,->] (-1.95,0.92) .. controls (-0.70,1.32) and (0.70,1.32) .. (1.95,0.92);
    \draw[blue,very thick,->] (1.95,-0.92) .. controls (0.70,-1.32) and (-0.70,-1.32) .. (-1.95,-0.92);
    \node[blue,font=\small] at (2.12,0.52) {$B$ 增强};
    \node[font=\small] at (0,-1.58) {铁芯磁化，磁场增强};
  \end{scope}
\end{tikzpicture}


\figcaption{电磁铁：通电螺线管插入软铁芯后，磁场被增强且可由电流控制}\label{fig:magnetic-material-electromagnet}
\end{minipage}
\end{center}

\begin{itemize}
  \item 电磁铁的强弱通常与电流大小、线圈匝数、铁芯材料有关。
  \item 软铁断电后磁性大多迅速减弱，适合做电磁铁；钢能较久保持磁性，适合做永久磁体。
  \item 高中阶段只作定性理解：铁芯不是“提供电流”，而是被磁化后增强原来的磁场。
\end{itemize}

\textbf{应用视角：} 电磁铁可用电流控制磁场强弱，是继电器、电磁起重机、扬声器、电磁阀和许多自动控制装置的基础。

\end{conceptbox}

\begin{conceptbox}{安培力}

一段通电直导线放在匀强磁场中会受到安培力。若磁场中的直导线长度为 $L$，电流为 $I$，磁感应强度为 $B$，电流方向与磁场方向夹角为 $\theta$，则
\[
\boxed{F = BIL\sin\theta}
\]

\begin{itemize}
  \item 当导线垂直于磁场时，$\theta = 90^\circ$，$F = BIL$。
  \item 当导线平行或反平行于磁场时，$\theta = 0^\circ$ 或 $180^\circ$，均有 $F = 0$。
  \item $L$ 指处在磁场中的直导线有效长度，不一定等于整根导线长度。
\end{itemize}

\textbf{适用边界：} 上式直接适用于匀强磁场中的通电直导线；若导线弯曲或磁场不均匀，应把导线分成许多足够短的小段，分别判断各段受力后再作矢量合成，不能把弯曲导线的弧长直接代入上式。

\begin{center}
% TikZ 图：安培力与夹角
% Day 21：磁场与安培力
\begin{tikzpicture}[scale=0.92,>=Latex]
  \coordinate (O) at (-0.5,-0.25);
  \foreach \y in {-1.05,-0.45,0.15,0.75} {
    \draw[->,blue] (-2.7,\y) -- (2.7,\y);
  }
  \node[blue] at (2.95,0.75) {$B$};
  \draw[very thick,red,->] (O) -- ++(2.25,1.55) node[above right] {$I$};
  \draw[dashed,gray] (O) -- ++(2.45,0);
  \draw[orange,thick,->] ($(O)+(0.82,0)$) arc (0:34.6:0.82);
  \node[orange] at ($(O)+(1.05,0.38)$) {$\theta$};

  \draw[purple,thick] (-2.05,1.12) circle (0.18);
  \draw[purple,thick] (-2.17,1.00) -- (-1.93,1.24);
  \draw[purple,thick] (-2.17,1.24) -- (-1.93,1.00);
  \node[purple,align=center] at (-1.25,1.13) {$F_A$\\向里};
  \node[font=\small] at (0,-1.65) {夹角 $\theta$ 是电流方向与磁场方向的夹角，$F_A=BIL\sin\theta$};
\end{tikzpicture}


\figcaption{安培力大小由电流方向与磁场方向的夹角决定}\label{fig:ampere-force-angle}
\end{center}

\textbf{左手定则：} 伸开左手，使磁感线穿入手心，四指指向电流方向，则拇指所指方向为安培力方向。

\begin{center}
% TikZ 图：左手定则示意
% Day 21：磁场与安培力
\begin{tikzpicture}[scale=0.9,>=Latex]
  \foreach \x in {-1.8,-1.1,-0.4,0.3,1.0,1.7} {
    \foreach \y in {-0.85,-0.15,0.55,1.25} {
      \draw[blue,thick] (\x-0.08,\y-0.08) -- (\x+0.08,\y+0.08);
      \draw[blue,thick] (\x-0.08,\y+0.08) -- (\x+0.08,\y-0.08);
    }
  }
  \node[blue] at (-1.35,1.60) {$B$ 向里};

  \draw[thick,fill=gray!12,rounded corners=8pt] (-0.85,-0.62) rectangle (0.30,0.58);
  \foreach \yy in {0.42,0.16,-0.10,-0.36} {
    \draw[thick,fill=gray!12,rounded corners=5pt] (0.22,\yy-0.10) rectangle (1.70,\yy+0.10);
  }
  \draw[thick,fill=gray!12,rounded corners=7pt] (-0.58,0.50) -- (-0.42,1.54) -- (-0.12,1.52) -- (-0.10,0.55) -- cycle;

  \draw[->,red,very thick] (0.35,0.42) -- (2.05,0.42) node[right] {$I$};
  \draw[->,purple,very thick] (-0.34,0.64) -- (-0.34,1.88) node[above] {$F_A$};
  \node[font=\small] at (0,-1.30) {四指指电流，拇指指安培力；磁场垂直纸面向里穿入手心};
\end{tikzpicture}


\figcaption{左手定则：手心迎磁场，四指指电流，拇指指安培力}\label{fig:left-hand-rule}
\end{center}

\textbf{方向判断清单：} 先看电流方向与磁场方向是否平行，若平行则 $F=0$；若不平行，再用左手定则判断方向。判断完要检查：安培力方向必须同时垂直于电流方向和磁场方向。

\end{conceptbox}

\begin{samepage}
\section*{四、公式与适用条件}

\textbf{方向判断前检查：}
\nopagebreak[4]
\begin{enumerate}
  \item 先明确磁感应强度 $B$ 的方向；磁感线某点的切线方向就是 $B$ 的方向。
  \item 通电导线受力要同时知道电流方向和磁场方向，再用左手定则。
  \item 代入安培力公式前先看电流与磁场夹角；垂直时用 $BIL$，不垂直时用 $BIL\sin\theta$。
  \item 判断电流产生的磁场时用右手螺旋定则，不要和左手定则混用。
  \item 图中“$\odot$”表示垂直纸面向外，“$\otimes$”表示垂直纸面向里。
\end{enumerate}
\end{samepage}

\begin{center}
\begin{tabular}{lll}
\toprule
公式 & 物理意义 & 适用条件 \\
\midrule
$B=\dfrac{F}{IL}$ & 磁感应强度定义 & 直导线、匀强场，且 $I\perp B$ \\
$F=BIL$ & 安培力最大值 & 直导线垂直匀强磁场 \\
$F=BIL\sin\theta$ & 安培力一般式 & 匀强磁场中的通电直导线 \\
$F=0$ & 无安培力 & $\theta=0^\circ$ 或 $180^\circ$ \\
\bottomrule
\end{tabular}
\end{center}

\begin{learningbox}{磁场题的方向对象分析清单}
磁场方向题先问“现在判断的是谁的方向”，再选择对应的手势规则：
\begin{enumerate}
  \item \textbf{判断电流产生的磁场：} 研究对象是“这根电流周围的 $B$”，用右手螺旋定则。
  \item \textbf{判断导线受到的安培力：} 研究对象是“这段有效导线”，先找外磁场 $B$ 和导线中电流 $I$，再用左手定则。
  \item \textbf{计算安培力大小：} 先确认有效长度 $L$ 是否全部在磁场中，再看电流与磁场夹角；垂直用 $F=BIL$，有夹角用 $F=BIL\sin\theta$。
\end{enumerate}
最后把安培力放回力学受力图中。若题目说导线平衡、悬浮、加速或沿导轨运动，后半段仍然是重力、弹力、摩擦力和牛顿定律的问题。

\textbf{来源理解：} $B=\dfrac{F}{IL}$ 是用通电导线受力来定义磁场强弱的操作式，要求直导线与匀强磁场垂直；$F=BIL\sin\theta$ 则把“不垂直时只有垂直于磁场的电流方向分量参与计算”的情况写进去。读题时先确认导线是否可看成直导线、所在区域是否可看成匀强磁场，再确认有效长度和夹角，最后把安培力作为一个真实力放入力学方程。
\end{learningbox}

\section*{五、典型模型}

\subsection*{模型 1：垂直纸面的磁场}

叉号“$\times$”像箭尾，表示磁场垂直纸面向里；点号“$\odot$”像箭头尖，表示磁场垂直纸面向外。遇到电磁题先把纸面内外方向翻译清楚。

\subsection*{模型 2：磁场也要按矢量叠加}

多个电流或磁体同时存在时，某点的磁场不能只凭距离远近粗略判断，而要把各个磁场看成矢量后合成。高中阶段常见处理顺序是：
\[
\text{分别判断方向}\rightarrow\text{比较大小}\rightarrow\text{矢量合成}.
\]
若两个磁场方向相同，合磁场增强；方向相反，合磁场减弱；方向成角度，则按平行四边形法则合成。

\textbf{图中条件：} 两根长直导线中的电流大小相等、方向相反，点 $P$ 到两导线距离相等。由右手螺旋定则，$P$ 点的 $\boldsymbol B_1$ 与 $\boldsymbol B_2$ 大小相等、方向关于竖直线对称，故合磁场沿竖直方向。

\begin{center}
% TikZ 图：磁场叠加
% Day 21：磁场与安培力
\begin{tikzpicture}[scale=0.9,>=Latex]
  \draw[thick] (-2.1,0) circle (0.20);
  \fill[red] (-2.1,0) circle (0.055);
  \node[font=\small] at (-3.05,-0.05) {$I_1$ 向外};
  \draw[blue,thick,->] (-2.1,0.65) arc (90:380:0.65);

  \draw[thick] (2.1,0) circle (0.20);
  \draw[red,thick] (1.98,-0.12) -- (2.22,0.12);
  \draw[red,thick] (1.98,0.12) -- (2.22,-0.12);
  \node[font=\small] at (3.05,-0.05) {$I_2$ 向里};
  \draw[blue,thick,->] (2.1,0.65) arc (90:-200:0.65);

  \fill[red] (0,1.15) circle (2pt) node[above=2pt] {$P$};
  \draw[->,blue,thick] (0,1.15) -- (-0.82,1.72) node[left] {$B_1$};
  \draw[->,blue,thick] (0,1.15) -- (0.82,1.72) node[right] {$B_2$};
  \draw[->,red,very thick] (0,1.15) -- (0,2.18) node[above] {$B_{\text{合}}$};

  \node[font=\small,align=center] at (0,-1.10)
    {$|I_1|=|I_2|$，且 $P$ 到两导线距离相等};
  \node[font=\small,align=center] at (0,-1.52)
    {先分别判断 $\boldsymbol B_1,\boldsymbol B_2$，再按平行四边形法则叠加};
\end{tikzpicture}


\figcaption{等大反向电流在对称点 $P$ 处的磁场叠加}\label{fig:magnetic-superposition}
\end{center}

\subsection*{模型 3：通电导线的受力平衡}

通电导线在磁场中常和重力、弹力、拉力一起平衡。为突出有效导线及其受力，以下示意图省略电源和其余供电导线，并忽略供电导线对研究对象的力。解题顺序：
\[
\text{定磁场方向} \to \text{定电流方向} \to \text{左手定则定安培力} \to \text{列平衡方程}
\]

\begin{center}
% TikZ 图：安培力与重力平衡
% Day 21：磁场与安培力
\begin{tikzpicture}[scale=0.92,>=Latex]
  \foreach \x in {-2.0,-1.2,-0.4,0.4,1.2,2.0} {
    \foreach \y in {-0.7,0.9} {
      \draw[blue,thick] (\x-0.09,\y-0.09) -- (\x+0.09,\y+0.09);
      \draw[blue,thick] (\x-0.09,\y+0.09) -- (\x+0.09,\y-0.09);
    }
  }
  \node[blue] at (2.55,1.05) {$B$ 向里};
  \draw[very thick,red,->] (-1.7,0.1) -- (1.7,0.1) node[right=2pt] {$I$};
  \draw[->,purple,thick] (0,0.1) -- (0,1.25) node[above] {$F_A$};
  \draw[->,black,thick] (0,0.1) -- (0,-1.05) node[below] {$mg$};
  \node at (0,-1.72) {安培力与重力平衡};
\end{tikzpicture}


\figcaption{水平导线在磁场中平衡：安培力可与重力平衡，供电回路未画出}\label{fig:ampere-balance}
\end{center}

\subsection*{模型 4：通电导体棒在导轨上运动}

导体棒中有电流，若处在磁场中，就受到安培力。安培力方向决定导体棒的加速度方向，是电动机模型的最简雏形。图中只画出磁场中的有效导体棒，电源和其余回路导线均省略。

\begin{center}
% TikZ 图：导体棒在导轨上受安培力
% Day 21：磁场与安培力
\begin{tikzpicture}[scale=0.92,>=Latex]
  \draw[thick] (-2,1) -- (2.2,1);
  \draw[thick] (-2,-1) -- (2.2,-1);
  \draw[line width=6pt,gray!35] (-0.2,-1) -- (-0.2,1);
  \draw[thick] (-0.2,-1) -- (-0.2,1);
  \draw[very thick,red,->] (-0.2,-0.9) -- (-0.2,0.9) node[near end,right=2pt] {$I$};
  \foreach \x in {-1.6,-0.8,0,0.8,1.6} {
    \foreach \y in {-0.55,0.25} {
      \draw[blue,thick] (\x-0.08,\y-0.08) -- (\x+0.08,\y+0.08);
      \draw[blue,thick] (\x-0.08,\y+0.08) -- (\x+0.08,\y-0.08);
    }
  }
  \node[blue] at (2.75,0.75) {$B$ 向里};
  \draw[->,purple,very thick] (-0.2,0) -- (-1.55,0) node[left] {$F_A$};
  \node[font=\small] at (0.55,1.32) {导体棒};
  \node at (0,-1.55) {导体棒受力后沿导轨运动};
\end{tikzpicture}


\figcaption{导体棒受安培力运动：磁场、电流、力三者两两垂直，完整供电回路未画出}\label{fig:rail-ampere-force}
\end{center}

\section*{六、例题精讲}

\begin{examplebox}{电流磁场的方向与叠加}
\textbf{题意分析：} 两根相互平行的长直导线垂直纸面放置，左侧导线电流垂直纸面向外，右侧导线电流垂直纸面向里，且电流大小相等。点 $P$ 位于两导线连线的垂直平分线上。判断两根导线在 $P$ 点产生的磁场方向，并确定合磁场方向。

\textbf{解题过程：} 分别对两根导线使用右手螺旋定则。左侧电流向外，磁场沿逆时针方向环绕，在 $P$ 点指向左上方；右侧电流向里，磁场沿顺时针方向环绕，在 $P$ 点指向右上方。由于电流大小相等且 $P$ 到两导线距离相等，两个分磁场大小相等，水平方向分量抵消，竖直方向分量相加，所以合磁场竖直向上。

\textbf{方法总结：} 多个电流共同产生磁场时，右手螺旋定则只负责判断每个分磁场方向；最后还要按平行四边形法则进行矢量合成。
\end{examplebox}

\begin{examplebox}{安培力大小与方向}
\textbf{题意分析：} 长度 $L=0.20\ \mathrm{m}$ 的水平直导线垂直放入匀强磁场，磁感应强度 $B=0.50\ \mathrm{T}$、方向垂直纸面向里，导线中电流 $I=3.0\ \mathrm{A}$、方向水平向右。求导线所受安培力的大小和方向。

\textbf{物理模型：} 直导线垂直匀强磁场，$F=BIL$。

\begin{center}
% TikZ 图：安培力例题示意
% Day 21：磁场与安培力
\begin{tikzpicture}[scale=0.82,>=Latex]
  \foreach \x in {-1.8,-1.1,-0.4,0.3,1.0,1.7} {
    \foreach \y in {-0.65,0.55} {
      \draw[blue,thick] (\x-0.08,\y-0.08) -- (\x+0.08,\y+0.08);
      \draw[blue,thick] (\x-0.08,\y+0.08) -- (\x+0.08,\y-0.08);
    }
  }
  \node[blue] at (2.35,0.72) {$B$ 向里};
  \draw[very thick,red,->] (-1.45,0) -- (1.45,0) node[right=2pt] {$I$};
  \draw[<->,black] (-1.45,-0.95) -- (1.45,-0.95) node[midway,below] {$L$};
  \draw[->,purple,very thick] (0,0) -- (0,1.08) node[above] {$F_A$};
  \node[font=\small] at (0,-1.42) {$I\perp B$，可直接用 $F_A=BIL$};
\end{tikzpicture}

\end{center}

\textbf{解题过程：}
\[
F = BIL = 0.50 \times 3.0 \times 0.20 = 0.30\ \mathrm{N}
\]
由左手定则，安培力方向竖直向上。

\textbf{方法总结：} 先判断夹角。若题中说“垂直放入”，可以直接用 $F=BIL$；若有夹角，必须乘 $\sin\theta$。
\end{examplebox}

\begin{examplebox}{通电导线的平衡}
\textbf{题意分析：} 一根水平导线长 $L=0.10\ \mathrm{m}$、质量 $m=0.020\ \mathrm{kg}$，处在垂直纸面向里的匀强磁场中。若希望导线静止悬浮，磁感应强度 $B=0.50\ \mathrm{T}$，求所需电流大小。取 $g=10\ \mathrm{m/s^2}$。

\textbf{物理模型：} 安培力与重力平衡。

\begin{center}
% TikZ 图：安培力平衡例题示意
% Day 21：磁场与安培力
\begin{tikzpicture}[scale=0.82,>=Latex]
  \foreach \x in {-1.8,-1.1,1.0,1.7} {
    \draw[blue,thick] (\x-0.08,0.67) -- (\x+0.08,0.83);
    \draw[blue,thick] (\x-0.08,0.83) -- (\x+0.08,0.67);
  }
  \foreach \x in {-1.8,-1.1,1.7} {
    \draw[blue,thick] (\x-0.08,-0.78) -- (\x+0.08,-0.62);
    \draw[blue,thick] (\x-0.08,-0.62) -- (\x+0.08,-0.78);
  }
  \node[blue] at (2.35,0.82) {$B$ 向里};
  \draw[very thick,red,->] (-1.45,0) -- (1.45,0) node[right=2pt] {$I$};
  \draw[->,purple,very thick] (0,0) -- (0,1.08) node[above] {$F_A$};
  \draw[->,black,thick] (0,0) -- (0,-0.90) node[pos=0.72,right=4pt] {$mg$};
  \node[font=\small] at (0,-1.45) {悬浮条件：$F_A=mg$};
\end{tikzpicture}

\end{center}

\textbf{解题过程：}
\[
F_A = mg,\qquad BIL = mg
\]
\[
I = \frac{mg}{BL}
= \frac{0.020\times 10}{0.50\times 0.10}
=4.0\ \mathrm{A}
\]

方向由左手定则确定：若磁场向里，安培力要向上，则电流应向右。

\textbf{方法总结：} 电磁平衡题仍然是力学题。安培力只是新增的一种力，先画受力图再列方程。
\end{examplebox}

\begin{examplebox}{导体棒在导轨上的加速}
\textbf{题意分析：} 光滑水平导轨间距 $L=0.30\ \mathrm{m}$，导体棒质量 $m=0.060\ \mathrm{kg}$，棒中电流 $I=4.0\ \mathrm{A}$，匀强磁场 $B=0.50\ \mathrm{T}$ 垂直纸面向里。若安培力水平向左，求导体棒加速度大小。

\textbf{物理模型：} 安培力提供水平方向合力，仍用牛顿第二定律。

\begin{center}
% TikZ 图：导体棒加速例题示意
% Day 21：磁场与安培力
\begin{tikzpicture}[scale=0.82,>=Latex]
  \draw[thick] (-2.0,0.9) -- (2.2,0.9);
  \draw[thick] (-2.0,-0.9) -- (2.2,-0.9);
  \foreach \x in {-1.55,-0.85,-0.15,0.55,1.25,1.95} {
    \foreach \y in {-0.45,0.45} {
      \draw[blue,thick] (\x-0.08,\y-0.08) -- (\x+0.08,\y+0.08);
      \draw[blue,thick] (\x-0.08,\y+0.08) -- (\x+0.08,\y-0.08);
    }
  }
  \node[blue] at (2.75,0.62) {$B$ 向里};
  \draw[very thick,red,->] (-0.2,-0.9) -- (-0.2,0.9) node[near end,right=2pt] {$I$};
  \draw[<->,black] (-1.88,-0.9) -- (-1.88,0.9) node[midway,left=2pt] {$L$};
  \draw[->,purple,very thick] (-0.2,0) -- (-1.35,0);
  \node[purple] at (-1.62,0.18) {$F_A$};
  \draw[->,orange,thick] (-0.2,-1.25) -- (-1.25,-1.25) node[midway,below] {$a$};
\end{tikzpicture}

\end{center}

\textbf{解题过程：}
导体棒所受安培力
\[
F_A=BIL=0.50\times4.0\times0.30=0.60\ \mathrm{N}.
\]
导轨光滑，水平方向合力就是安培力：
\[
a=\frac{F_A}{m}=\frac{0.60}{0.060}=10\ \mathrm{m/s^2}.
\]

\textbf{方法总结：} “电磁题”并不意味着力学消失。判断出安培力方向和大小后，后面就是受力分析和牛顿第二定律。
\end{examplebox}

\section*{七、变式训练}

\begin{enumerate}
  \item 长 $0.30\ \mathrm{m}$ 的导线中通有 $2.0\ \mathrm{A}$ 电流，垂直放入 $0.40\ \mathrm{T}$ 的匀强磁场中，求安培力大小。
  \item 若上题中导线与磁场夹角为 $30^\circ$，安培力变为多少？
  \item 一根质量 $0.05\ \mathrm{kg}$、长 $0.25\ \mathrm{m}$ 的水平导线悬挂在匀强磁场中，若电流为 $4\ \mathrm{A}$ 时恰好平衡，求磁感应强度大小。
  \begin{center}
  % TikZ 图：安培力与重力平衡
% Day 21：磁场与安培力
\begin{tikzpicture}[scale=0.92,>=Latex]
  \foreach \x in {-2.0,-1.2,-0.4,0.4,1.2,2.0} {
    \foreach \y in {-0.7,0.9} {
      \draw[blue,thick] (\x-0.09,\y-0.09) -- (\x+0.09,\y+0.09);
      \draw[blue,thick] (\x-0.09,\y+0.09) -- (\x+0.09,\y-0.09);
    }
  }
  \node[blue] at (2.55,1.05) {$B$ 向里};
  \draw[very thick,red,->] (-1.7,0.1) -- (1.7,0.1) node[right=2pt] {$I$};
  \draw[->,purple,thick] (0,0.1) -- (0,1.25) node[above] {$F_A$};
  \draw[->,black,thick] (0,0.1) -- (0,-1.05) node[below] {$mg$};
  \node at (0,-1.72) {安培力与重力平衡};
\end{tikzpicture}

  \end{center}
  \item 判断：磁场垂直纸面向外，导线中电流向右，则安培力方向是向上还是向下？
  \begin{center}
  \begin{tikzpicture}[scale=0.85,>=Stealth]
  \foreach \x in {0,0.6,1.2,1.8,2.4,3.0}
    \foreach \y in {0,1.2}
      \node[blue] at (\x,\y) {$\odot$};
  \draw[red,very thick,->] (-0.15,0.6) -- (3.25,0.6) node[right] {$I$};
  \draw[->,black,thick] (1.55,0.48) -- (1.55,-0.45) node[below] {$F$};
  \node[font=\small] at (1.5,1.62) {磁场向外，电流向右};
\end{tikzpicture}

  \end{center}
  \item 通电螺线管中插入软铁芯后，磁场强弱如何变化？断电后为什么常选软铁制作电磁铁铁芯？
  \item 从某一环形电流的正面观察，电流沿逆时针方向流动。环心处的磁场指向纸内还是纸外？这一面相当于 $N$ 极还是 $S$ 极？
  \item 从通电螺线管左端观察，线圈中的电流沿顺时针方向流动。螺线管左端是什么磁极？管内磁场沿轴线指向左端还是右端？
\end{enumerate}

\begin{learningbox}{变式训练参考结果}
\begin{enumerate}[itemsep=0.45em,topsep=0.3em,parsep=0pt]
  \item $F=BIL=0.24\ \mathrm N$。
  \item $F=BIL\sin30^\circ=0.12\ \mathrm N$。
  \item 平衡时 $BIL=mg$，故 $B=0.50\ \mathrm T$；安培力方向竖直向上。
  \item 由左手定则（或 $\boldsymbol F=I\boldsymbol L\times\boldsymbol B$）判断，安培力方向向下。
  \item 插入软铁芯后磁场显著增强。软铁磁导率较大且剩磁较小，通电时易被磁化，断电后又较易退磁，因此适合作电磁铁铁芯。
  \item 环心处磁场垂直纸面向外；从这一面看，磁感线从线圈穿出，因此这一面相当于 $N$ 极。
  \item 左端为 $S$ 极；管内磁场由 $S$ 极指向 $N$ 极，即沿轴线指向右端。
\end{enumerate}
\end{learningbox}

\section*{八、竞赛入门补充}

\begin{competitionbox}
\textbf{从安培力看电动机的基本思想}

电动机的基本物理图像是：通电线圈在磁场中受到安培力产生的转矩。每一段载流导线受到的是沿导线分布的磁场力，通常用其合力表示。线圈两条相对的有效边中电流方向相反，所受安培力大小相等、方向相反，但作用线不重合，因而构成\textbf{力偶}：线圈所受合力可以为零，转矩却不为零。

\begin{center}
% TikZ 图：通电线圈受力矩
% Day 21：磁场与安培力
\begin{tikzpicture}[scale=1.02,>=Latex,line join=round]
  \path[use as bounding box] (-3.55,-1.75) rectangle (3.55,1.70);

  \foreach \y in {-1.05,-0.35,0.35,1.05} {
    \draw[->,blue,thick] (-3.25,\y) -- (3.25,\y);
  }
  \node[blue] at (3.45,1.05) {$B$};

  \coordinate (A) at (-1.25,-0.88);
  \coordinate (B) at (1.05,-0.88);
  \coordinate (C) at (1.55,0.78);
  \coordinate (D) at (-0.75,0.78);
  \coordinate (E) at (-0.70,-0.58);
  \coordinate (F) at (1.60,-0.58);
  \coordinate (G) at (2.10,1.08);
  \coordinate (H) at (-0.20,1.08);

  \draw[thick,gray!65,dashed] (E) -- (F) -- (G) -- (H) -- cycle;
  \draw[thick,fill=gray!8] (A) -- (B) -- (C) -- (D) -- cycle;
  \draw[thick] (A) -- (E) (B) -- (F) (C) -- (G) (D) -- (H);

  \draw[thick,dashed] (0.10,-1.28) -- (0.42,1.40);

  \draw[->,red,very thick] (A) -- (D);
  \draw[->,red,very thick] (G) -- (F);

  \draw[purple,thick] (-1.62,0.03) circle (0.18);
  \draw[purple,thick] (-1.74,-0.09) -- (-1.50,0.15);
  \draw[purple,thick] (-1.74,0.15) -- (-1.50,-0.09);
  \node[purple,align=center] at (-2.42,0.03) {$F_A$\\向里};

  \draw[purple,thick] (2.02,0.03) circle (0.18);
  \fill[purple] (2.02,0.03) circle (0.045);
  \node[purple,align=center] at (2.80,0.03) {$F_A$\\向外};

  \draw[->,orange,thick] (0.78,1.14) arc (105:-220:0.46);
  \node[orange] at (1.32,1.38) {转动};
  \node[font=\small] at (0,-1.58) {线圈两侧受一里一外的安培力，形成力矩};
\end{tikzpicture}


\figcaption{通电线圈两侧的安培力构成力偶并产生转矩}\label{fig:motor-coil-torque}
\end{center}

\textbf{转矩的定量表达：} 设矩形线圈有 $N$ 匝，单匝面积为 $S$，电流为 $I$，线圈平面的法线与磁场方向夹角为 $\alpha$。两条有效边所受安培力形成力偶，可得
\[
\boxed{M=NBIS\sin\alpha}.
\]
当线圈平面与磁场垂直时，法线与 $\boldsymbol B$ 平行，$\alpha=0$，转矩为零；当线圈平面与磁场平行时，$\alpha=90^\circ$，转矩最大。这里的 $M$ 是转矩大小，单位为 $\mathrm{N\cdot m}$。

\textbf{微观来源：} 安培力可以理解为载流导线中许多运动电荷所受磁场力的宏观合力。下一节学习洛伦兹力后，会把“单个运动电荷”和“整段电流”两个层次联系起来。

\textbf{例：} $N=20$ 匝的矩形线圈面积 $S=6.0\times10^{-3}\ \mathrm{m^2}$，通有 $I=2.0\ \mathrm A$ 的电流，置于 $B=0.50\ \mathrm T$ 的匀强磁场中。若线圈法线与磁场夹角 $\alpha=30^\circ$，则
\[
M=NBIS\sin\alpha
=20\times0.50\times2.0\times6.0\times10^{-3}\times\sin30^\circ
=6.0\times10^{-2}\ \mathrm{N\cdot m}.
\]

\textbf{进一步理解：} 线圈转动时 $\alpha$ 不断变化，转矩也随之变化。真实电动机需要换向器或电子控制，使线圈越过转矩为零的位置后，电流适时换向，从而保持连续转动。
\end{competitionbox}

\textbf{拓展检查：} 上式中的 $\alpha$ 是\textbf{线圈平面法线}与磁场的夹角，不是线圈平面与磁场的夹角。若题目给出线圈平面与磁场夹角 $\beta$，则 $\alpha=90^\circ-\beta$，应写成 $M=NBIS\cos\beta$。

\section*{九、易错点}

\begin{mistakebox}
\begin{enumerate}
  \item \textbf{把磁感线当成真实存在的线：} 读图时只做两件事：沿箭头读某点的磁场方向，比较疏密读强弱趋势；不要把它画成电流、磁体或粒子的实际运动轨迹，更不要说“电流沿磁感线运动”。
  \item \textbf{混用左右手定则：} 先判断题目问“电流产生磁场”还是“磁场对电流的力”；前者用右手，后者用左手。
  \item \textbf{忘记夹角：} 写 $F=BIL$ 前先确认电流与磁场垂直；一般式是 $F=BIL\sin\theta$。
  \item \textbf{把 $L$ 当成整根导线长度：} 先圈出导线在磁场中的那一段，$L$ 是有效长度。
  \item \textbf{忽略安培力方向：} 当电流方向与磁场方向不平行、安培力不为零时，安培力方向同时垂直于电流方向和磁场方向；若二者平行，则 $F=0$，不再讨论力的方向。
  \item \textbf{把地理北极当成地磁 $N$ 极：} 指南针 $N$ 极指向地理北方，说明地理北极附近在磁性上相当于地磁 $S$ 极。
  \item \textbf{误以为铁芯本身产生电流：} 铁芯被磁化后增强螺线管磁场；判断电磁铁强弱时仍先看线圈电流和匝数。
\end{enumerate}
\end{mistakebox}

\section*{十、今日小结}

\begin{enumerate}
  \item 判断磁场方向时，先看小磁针静止时 $N$ 极指向；磁感线切线方向就是该点磁场方向。
  \item 地磁场可近似看成条形磁体磁场，地理北极附近相当于地磁 $S$ 极；读图时要分清地理南北和地磁南北。
  \item 磁感应强度 $B$ 描述磁场强弱和方向，单位是 $\mathrm{T}$；列式前先判断 $B$ 与电流方向是否垂直。
  \item 电流产生磁场用右手螺旋定则，先确定研究的是直导线、环形电流还是螺线管。
  \item 铁磁性材料容易被磁化；通电螺线管插入软铁芯可形成电磁铁。
  \item 通电导线在磁场中受安培力：$F=BIL\sin\theta$；若导线与磁场平行，安培力为零。
  \item 判断安培力方向用左手定则，并把安培力作为真实外力纳入力学受力分析。
\end{enumerate}

\textbf{一分钟复盘：} 任选一根通电导线，先在图上标出电流 $I$ 和外磁场 $B$，用左手判断安培力 $F$；若题目问“电流产生的磁场”，再改用右手螺旋定则判断磁感线方向。复盘时要说清：现在判断的是“受力”，还是“产生磁场”。
